\section{Security Test Plan}
\label{sec:testplan}

This section plans the security tests of the SecureBank~PT release
candidate. The scope is the product backend and its clients: the mobile
application (P1), the back-office desktop (P2), the \ac{API} gateway (P3),
the \ac{IAM} service (P4), the banking services (P5) and the \ac{CI/CD}
pipeline (P6) of Figure~\ref{fig:dfd}. The core banking system and SIBS
(EE3) belong to third parties and are replaced by a simulator. All tests
run in a staging copy of TB3 with its own \ac{HSM} partition and
synthetic customers, so no real personal data is used. A release enters
testing when the pipeline gates of Section~\ref{sec:lifecycle} pass, and
it leaves testing when every requirement has a passed test and no
critical or high finding is open (\mbox{REQ-23}). The plan has three
parts: a penetration-testing plan, a fuzz-testing plan and the mapping of
every requirement of Table~\ref{tab:req} to at least one test. Test cases
follow the ISO/IEC/IEEE 29119-3 template~\cite{iso29119}, and
Annex~\ref{ann:tc04} gives one test case in full.

\subsection{Penetration-Testing Plan}

\textbf{Strategy.} The penetration test is grey-box and threat-led. The
testers receive the \ac{DFD}, the threat table of Annex~B, the OpenAPI
description of the gateway, release builds of the Android and iOS
applications, and test accounts for two customers and two Admins. They do
not receive the source code, so the test shows what an attacker with an
account and the public application can do. The engagement follows the
phases of \ac{PTES}~\cite{ptes}: pre-engagement, reconnaissance, threat
modelling, vulnerability analysis, exploitation, post-exploitation and
reporting. The OWASP \ac{WSTG}~\cite{wstg} is the checklist for the
\ac{API} and the back-office, and the \ac{MASTG}~\cite{mastg} for the
mobile application. Every finding is reported with reproducible steps,
the threat it confirms, its impact and a proposed fix.

\textbf{Tools.} nmap and testssl.sh map the exposed services and the TLS
configuration. Burp Suite Professional is the main intercepting proxy,
with the Autorize extension for access-control tests, and OWASP ZAP runs
the automated scans that the pipeline already uses. jwt\_tool tests the
tokens. On the mobile side, jadx and MobSF are used for static review of
the release builds, apktool for repackaging, and Frida with objection for
runtime instrumentation on a rooted Android device and a jailbroken iOS
device. sqlmap is used only on the back-office search fields, and k6
generates controlled load for the rate-limit tests.

\textbf{Attack types.} Table~\ref{tab:pentest} lists five attack types,
each derived from the threats of Annex~B, with the components most
exposed to it. The gateway and the banking services (P3, P5) are the
most sensitive components overall, since any enrolled device can reach
them. P4 comes next, because one flaw there gives account takeover.
The mobile application (P1) runs on a device the attacker may control,
so its protections can only delay an attacker, and the backend must
not trust any value that the application computes.

\begin{table*}[!t]
\caption{Penetration-testing plan: attack types, sensitive components, tools and pass criteria}
\label{tab:pentest}
\centering
\scriptsize
\begin{tabularx}{\textwidth}{@{}L{0.45cm}L{2.6cm}L{1.2cm}L{3.1cm}L{2.7cm}Y@{}}
\toprule
\textbf{ID} & \textbf{Attack type} & \textbf{Threats} & \textbf{Most sensitive components} & \textbf{Tools} & \textbf{Pass criterion (requirements)} \\
\midrule
A1 & Broken object-level authorisation and role confusion: change account, card or transaction identifiers; call back-office endpoints with a customer token & T14, T19 & P5 account, card and statement endpoints; P3 routing; back-office endpoints of P2 & Burp with Autorize, Python scripts with two customers & Every request for an object of another customer returns 403 or 404 with no data; no role reaches the endpoints of the other role (REQ-04, REQ-17) \\
A2 & Authentication and session attacks: \ac{PIN} guessing, credential stuffing, token theft and replay, step-up bypass & T01, T02, T04, T17 & P4 \ac{IAM}; F3 token handling; P1 login flow & Burp Intruder at a limited rate, jwt\_tool, Frida & Lockout after five wrong \acp{PIN}; tokens older than 5~minutes or without proof of possession rejected; one factor never enough (REQ-13, REQ-15, REQ-16, REQ-21) \\
A3 & Transaction tampering and replay: change amount or payee after signing, replay signed orders, reuse nonces & T06, T07, T10 & P1 signing code; F3; P5 transfer and MB~WAY services & Frida hooks on the signing call, Burp Repeater & Altered or replayed orders return 401 or 409 and leave balances unchanged (REQ-06, REQ-07) \\
A4 & Channel interception and client hardening: user-installed CA, pinning bypass, repackaged application, data left on the device & T03, T12, T13 & P1; DS1 keystore; F3 & mitmproxy, objection, apktool, MobSF & No request is readable through the proxy without a pinning bypass; no token, key or account data outside the keystore; a repackaged build cannot enrol (REQ-01, REQ-03, REQ-15) \\
A5 & Back-office abuse: freeze an account or open \ac{BtG} without a second Admin, approve one's own request, alter audit records, inject into search fields & T11, T20 & P2; P5 administrative services; DS3 audit log & Burp, sqlmap, direct database access granted for the test & Single-Admin actions refused; tampering breaks the hash chain; no injection reaches the database (REQ-09, REQ-18) \\
\bottomrule
\end{tabularx}
\end{table*}

\textbf{Staying within the authorised scope.} No test starts before the
CISO and the PO sign the rules of engagement. This document names the
testers, their source IP addresses, the time window, the staging host
names and the hashes of the builds under test. It also lists what is out
of scope: production, the core banking system and SIBS, the app-store
accounts, real customers and staff (no social engineering in this
engagement), and load above twice the \mbox{REQ-12} threshold. The
gateway accepts test traffic only from the listed addresses, and the
\ac{SIEM} team is told the window in advance, except for the
\mbox{REQ-21} alert drill. Testers stop and report within one hour if a
test reaches a system that is not on the list, exposes data that looks
real, or degrades shared infrastructure. Critical findings go to the
\ac{PSIRT} within 24~hours, and evidence is stored encrypted and deleted
90~days after the final report. Under \ac{DORA} Art.~26, a bank that
its competent authority identifies for it must also run threat-led
penetration testing on live production systems at least every three
years, following TIBER-EU~\cite{tibereu}. That test has its own control
team and is outside this plan.

\subsection{Fuzz-Testing Plan}

\textbf{Strategy and tools.} Fuzzing is used at two levels, and both are
coverage-guided or schema-guided rather than purely random~\cite{miller1990}.
Inside the backend, Jazzer~\cite{jazzer} (Table~\ref{tab:lifecycle})
fuzzes the JVM parsers and validators of orders, tokens and signatures
in-process. Each target receives bytes from the fuzzer and calls one
parser, and the seeds are valid requests captured from the functional
test suite. Through the gateway, Schemathesis~\cite{schemathesis}
generates requests from the OpenAPI description and chains them
(login, token, transfer), so that the fuzzer reaches the transfer
endpoints with a valid session. Native libraries of the mobile
application, such as the QR-code decoder, are fuzzed with
AFL++~\cite{aflpp2020} under AddressSanitizer, as in
Section~\ref{sec:fuzzrun}.

\textbf{Monitors.} A fuzzer finds crashes and memory faults, not logic
flaws, so each campaign adds oracles that check the behaviour of the
system: an uncaught exception or a sanitizer report; an HTTP~5xx or a
stack trace in a response; a response slower than 2~seconds; a change
in the total balance of the test accounts without an accepted signed
order; data of another test customer in a response; and a break in the
audit-log chain. Jazzer runs nightly for 8~hours per target, Schemathesis
for 30~minutes on each staging build, and each crash becomes a
regression test once it is fixed. Table~\ref{tab:fuzz} lists the ten
attack vectors and the expected behaviour of the system for each one.

\begin{table*}[!t]
\caption{Fuzz-testing plan: attack vectors, example inputs and expected behaviour}
\label{tab:fuzz}
\centering
\scriptsize
\begin{tabularx}{\textwidth}{@{}L{0.8cm}L{2.6cm}L{5.0cm}YL{1.3cm}@{}}
\toprule
\textbf{ID} & \textbf{Vector (target)} & \textbf{Example inputs} & \textbf{Expected behaviour} & \textbf{REQ} \\
\midrule
FZ-01 & Amount field of transfers and MB~WAY payments (P5 order parser) & \texttt{-0.01}; \texttt{0.001}; \texttt{1e309}; \texttt{NaN}; \texttt{99999999999999999999.99}; \texttt{"12,50"} & 400 with a generic message; only positive decimals with two places and within the account limit are accepted; no debit and no 5xx & REQ-06, REQ-08 \\
FZ-02 & Duplicate and conflicting JSON keys (P3 and P5 parsers) & \texttt{\{"amount":"10.00", "amount":"9000.00"\}}; the same key with different case & 400 at the gateway; the service never executes a value other than the one covered by the signature & REQ-06 \\
FZ-03 & Nesting depth and body size (P3) & 10{,}000 nested arrays \texttt{[[[[...]]]]}; a 10~MB body; a 1~MB string value & 400 or 413 within 2~seconds; no stack overflow; other sessions keep their response times & REQ-11, REQ-12 \\
FZ-04 & IBAN and payee name (P5 verification of payee) & \texttt{PT50} followed by 1{,}000 digits; IBAN with a NUL byte; name with U+202E (right-to-left override); Cyrillic homoglyphs in a known name & Malformed IBAN rejected with 400; names compared after Unicode normalisation; a non-matching name triggers the second confirmation & REQ-20 \\
FZ-05 & Signature field (P5 signature check) & empty value; truncated DER; 1~MB of base64; a valid signature of another order; $r=0$ or $s=0$ & 401; no transfer; no exception reaches the logs as a stack trace & REQ-06 \\
FZ-06 & Nonce, idempotency key and timestamp (P5) & reused nonce; \texttt{ts=0}; \texttt{ts=-1}; \texttt{ts=9223372036854775807}; \texttt{2026-12-31T23:59:60Z}; 10~kB idempotency key & 409 on reuse; 400 on malformed values; any timestamp older than 60~seconds or in the future rejected & REQ-07 \\
FZ-07 & Access token and proof of possession (P4) & header with \texttt{"alg":"none"}; \texttt{"kid":"../../dev/null"}; 64~kB token; token expired by one second; missing or reused proof & 401 for all; no 5xx; no token accepted without a valid proof from the enrolled device & REQ-16, REQ-17 \\
FZ-08 & Object identifiers in paths and bodies (P3, P5) & \texttt{0}; \texttt{-1}; \texttt{9223372036854775808}; \texttt{../}; \texttt{\%00}; identifier of the second test customer & 403 or 404; the oracle finds no IBAN or name of the other customer in any response & REQ-04 \\
FZ-09 & \ac{PIN} entry and login (P4) & empty; five or seven digits; Arabic-Indic digits U+0661--U+0666; \texttt{000000} followed by a newline; 10~kB string & 400 or 401; every attempt counts towards the five-attempt lockout, whatever its encoding & REQ-13, REQ-15 \\
FZ-10 & Deep links and QR payloads (P1 and its native decoder) & \texttt{securebank://transfer?}\allowbreak\texttt{amount=1e9\&auto=1}; QR with a TLV length of \texttt{0xFF} past the end; MB~WAY number \texttt{+351} followed by 4{,}000 digits & The application opens a pre-filled form that still needs biometrics or the \ac{PIN}; malformed codes rejected; no crash and no sanitizer report & REQ-06, REQ-15 \\
\bottomrule
\end{tabularx}
\end{table*}

\subsection{Requirement-to-Test Mapping}

Table~\ref{tab:tests} maps each of the 23 requirements to at least one
test with a unique identifier. Twenty-one test cases cover the catalogue.
Three of them cover more than one requirement: TC-05 checks the three
customer duties, TC-18 checks both the four-eyes rule and its audit
records, and TC-21, the independent penetration test, re-runs attack
types A1 to A5 and so repeats the checks of \mbox{REQ-01}, \mbox{REQ-04},
\mbox{REQ-06} and \mbox{REQ-17}. The \emph{Link} column names the attack
types and fuzz vectors that feed each test.

\begin{table*}[!t]
\caption{Requirement-to-test mapping (every requirement of Table~\ref{tab:req} has at least one test)}
\label{tab:tests}
\centering
\scriptsize
\begin{tabularx}{\textwidth}{@{}L{0.75cm}YL{2.3cm}L{1.35cm}L{1.35cm}L{1.9cm}@{}}
\toprule
\textbf{Test} & \textbf{Action and pass criterion} & \textbf{Technique} & \textbf{Link} & \textbf{When} & \textbf{Requirements} \\
\midrule
TC-01 & Route the application through a proxy with a user-installed CA, and through a server that offers only TLS~1.2: the application aborts both connections and sends no request. & Pentest & A4 & Each release & REQ-01 \\
TC-02 & Restore a backup on an isolated host without \ac{HSM} access and read the database files: no account or personal field is readable; configuration shows AES-256-GCM with keys in the \ac{HSM}. & Inspection, review & -- & Yearly & REQ-02 \\
TC-03 & Log in and make a transfer on a rooted device, then search the sandbox, logs and backups: no token, key or account data outside the keystore. & Pentest (MobSF, objection) & A4 & Each release & REQ-03 \\
TC-04 & Customer A requests every account, card and transaction identifier of customer B on every endpoint: 403 or 404 and no data of B (Annex~\ref{ann:tc04}). & Pentest, fuzzing & A1, FZ-08 & Nightly (scripted) & REQ-04 \\
TC-05 & Enrol a test customer and check the stored record: acceptance of the three duties, contract version and timestamp are present and cannot be skipped. & Functional, review & -- & Each release & REQ-05, REQ-19, REQ-22 \\
TC-06 & Change amount, payee, timestamp or nonce after signing: 401 and balances unchanged. & Pentest, fuzzing & A3, FZ-01, FZ-02, FZ-05 & Nightly & REQ-06 \\
TC-07 & Resend a signed order with the same nonce and idempotency key, and one signed 61~seconds earlier: 409 and 400; one debit only. & Pentest, boundary & A3, FZ-06 & Nightly & REQ-07 \\
TC-08 & Stop P5 or drop F11 at each step of a transfer: no partial debit or credit; reconciliation finds no difference. & Fault injection & -- & Each release & REQ-08 \\
TC-09 & Generate 1{,}000 events, then edit one record and delete another in DS3: chain verification fails exactly at those records; every record has actor, time, operation and outcome. & Functional, tampering & A5 & Each release & REQ-09 \\
TC-10 & Deploy an unsigned artefact and one whose \ac{SBOM} signature does not match: no node accepts either. & Functional & -- & Each pipeline change & REQ-10 \\
TC-11 & k6 with 1{,}000 sessions at 100 requests per second for 30~minutes: at least 99.9\% of valid requests complete within 2~seconds. & Load & FZ-03 & Each release & REQ-11 \\
TC-12 & One device sends 20 requests per second for 60~seconds while 50 devices send normal traffic: the flooding device receives 429 above 10 per second; the others are not affected. & Load, boundary & -- & Each release & REQ-12 \\
TC-13 & Enter four, then five wrong \acp{PIN}: the fifth blocks \ac{PIN} entry on that device; entry is still blocked at 14:59 and allowed at 15:00; matrix-card recovery works during the block. & Boundary & A2, FZ-09 & Nightly & REQ-13 \\
TC-14 & Stop the primary data centre during a k6 run: transfers resume from the second region within 30~minutes; every committed transaction is present in both ledgers. & Failover drill & -- & Twice a year & REQ-14 \\
TC-15 & Log in and transfer with the device key alone, the \ac{PIN} alone from a non-enrolled device, and biometrics without the key: all rejected; key with \ac{PIN} accepted. & Functional (negative), pentest & A2, FZ-10 & Nightly & REQ-15 \\
TC-16 & Present a token at 4:59 and at 5:01 minutes, and from a second device without the key: accepted, rejected, rejected. & Boundary, pentest & A2, FZ-07 & Nightly & REQ-16 \\
TC-17 & Call every back-office endpoint with a customer token and every payment endpoint with an Admin token: 403 for all. & Pentest (role matrix) & A1 & Nightly & REQ-17 \\
TC-18 & An Admin freezes an account and opens \ac{BtG} alone, then approves their own request: refused; with a second Admin: accepted and logged with both identities. & Functional, pentest & A5 & Each release & REQ-18, REQ-09 \\
TC-19 & Pay a new payee whose name does not match the verification of payee, including homoglyph names: second confirmation shown; a matching name shows none. & UI test & FZ-04 & Each release & REQ-20 \\
TC-20 & Replay 25 failed logins in one minute from one source, then 20: the \ac{SIEM} alerts within 60~seconds for the first and not for the second. & Alert drill & A2 & Quarterly & REQ-21 \\
TC-21 & Independent grey-box penetration test of A1 to A5 on the release candidate: no open critical or high finding, or residual risk signed by the PO. & Pentest & A1--A5 & Each major release & REQ-23, REQ-01, REQ-04, REQ-06, REQ-17 \\
\bottomrule
\end{tabularx}
\end{table*}

\textbf{Coverage and limits.} The tests marked \emph{nightly} run in the
pipeline against staging, so a regression in access control, signatures
or tokens is caught within a day. Some requirements cannot be fully
tested. The customer duties (\mbox{REQ-05}, \mbox{REQ-19},
\mbox{REQ-22}) are checked only through the enrolment records, since the
plan cannot observe what a customer does outside the application.
Ransomware (T18) is covered by the failover drill but not simulated with
real malware, and staging uses a smaller \ac{HSM} and network than
production, so TC-11 must be repeated on the production hardware before
peak periods. Fuzzing covers the parsers well but finds an authorisation
flaw only when an oracle checks for it, which is why FZ-08 and TC-04 share
the same two test customers.
